
%%
%% Part 3 – Practical Implementation of Diffusion Models
%%
%% In this chapter we translate the theoretical insights from Part~2
%% into concrete code.  We describe the U‑Net backbone used for
%% denoising, implement the forward diffusion and reverse sampling
%% procedures, and provide training routines.  Examples and tables
%% illustrate parameter choices and visual results.

\section{Implementation of Diffusion Models}

This chapter details the practical steps required to implement
Denoising Diffusion Probabilistic Models (DDPM) and their deterministic
counterpart, Denoising Diffusion Implicit Models (DDIM).  We begin
with the neural architecture, proceed to the variance schedule and
forward perturbation, outline the training loop, describe the sampling
routines and conclude with visual results and remarks on parameter
settings.

\subsection{U‑Net Architecture for Noise Prediction}

The U‑Net serves as the core network for predicting noise in the
reverse process.  It consists of an encoder and decoder connected via
skip connections.  Each level contains two residual blocks followed by
an attention mechanism, allowing the network to capture both local and
non‑local features.  Time embeddings are injected into every block to
condition the network on the current timestep.

\begin{lstlisting}[language=Python, caption=U‑Net Class Definition]
class UNet(nn.Module):
    def __init__(self, in_channels=3, out_channels=3, 
                 time_emb_dim=256, base_channels=64):
        super().__init__()
        # Time embedding
        self.time_emb = nn.Sequential(
            SinusoidalTimeEmbedding(time_emb_dim),
            nn.Linear(time_emb_dim, time_emb_dim),
            nn.SiLU(),
            nn.Linear(time_emb_dim, time_emb_dim)
        )
        # Encoder
        self.enc1 = ResBlock(in_channels, base_channels, time_emb_dim)
        self.enc2 = ResBlock(base_channels, base_channels*2, time_emb_dim)
        self.enc3 = ResBlock(base_channels*2, base_channels*4, time_emb_dim)
        self.enc4 = ResBlock(base_channels*4, base_channels*8, time_emb_dim)
        # Decoder
        self.dec1 = ResBlock(base_channels*8 + base_channels*4, base_channels*4, time_emb_dim)
        self.dec2 = ResBlock(base_channels*4 + base_channels*2, base_channels*2, time_emb_dim)
        self.dec3 = ResBlock(base_channels*2 + base_channels, base_channels, time_emb_dim)
        self.dec4 = ResBlock(base_channels + in_channels, base_channels, time_emb_dim)
        # Down/Up sampling layers
        self.down1 = nn.Conv2d(base_channels, base_channels, 4, 2, 1)
        self.down2 = nn.Conv2d(base_channels*2, base_channels*2, 4, 2, 1)
        self.down3 = nn.Conv2d(base_channels*4, base_channels*4, 4, 2, 1)
        self.up1 = nn.ConvTranspose2d(base_channels*8, base_channels*4, 4, 2, 1)
        self.up2 = nn.ConvTranspose2d(base_channels*4, base_channels*2, 4, 2, 1)
        self.up3 = nn.ConvTranspose2d(base_channels*2, base_channels, 4, 2, 1)
        # Output layer
        self.out = nn.Conv2d(base_channels, out_channels, 3, 1, 1)

    def forward(self, x, t):
        t_emb = self.time_emb(t)
        e1 = self.enc1(x, t_emb)
        e2 = self.enc2(self.down1(e1), t_emb)
        e3 = self.enc3(self.down2(e2), t_emb)
        e4 = self.enc4(self.down3(e3), t_emb)
        d1 = self.dec1(torch.cat([self.up1(e4), e3], dim=1), t_emb)
        d2 = self.dec2(torch.cat([self.up2(d1), e2], dim=1), t_emb)
        d3 = self.dec3(torch.cat([self.up3(d2), e1], dim=1), t_emb)
        d4 = self.dec4(torch.cat([d3, x], dim=1), t_emb)
        return self.out(d4)
\end{lstlisting}

\paragraph{Residual Blocks.}  Each \code{ResBlock} consists of group
normalisation, SiLU activation and two convolutional layers.  A time
embedding is added to the first convolution output to condition the
feature map on the timestep.

\begin{lstlisting}[language=Python, caption=Residual Block]
class ResBlock(nn.Module):
    def __init__(self, in_channels, out_channels, time_emb_dim):
        super().__init__()
        self.norm1 = nn.GroupNorm(32, in_channels)
        self.conv1 = nn.Conv2d(in_channels, out_channels, 3, 1, 1)
        self.norm2 = nn.GroupNorm(32, out_channels)
        self.conv2 = nn.Conv2d(out_channels, out_channels, 3, 1, 1)
        self.time_proj = nn.Linear(time_emb_dim, out_channels)
        self.residual = nn.Conv2d(in_channels, out_channels, 1) if in_channels != out_channels else nn.Identity()

    def forward(self, x, t_emb):
        h = F.silu(self.norm1(x))
        h = self.conv1(h)
        h = h + self.time_proj(t_emb)[..., None, None]
        h = F.silu(self.norm2(h))
        h = self.conv2(h)
        return h + self.residual(x)
\end{lstlisting}

\paragraph{Time Embedding.}  A sinusoidal embedding encodes the
position of the timestep into a high‑dimensional vector.  This
embedding is passed through two linear layers to produce a learned
embedding.

\begin{lstlisting}[language=Python, caption=Sinusoidal Time Embedding]
class SinusoidalTimeEmbedding(nn.Module):
    def __init__(self, dim):
        super().__init__()
        self.dim = dim

    def forward(self, t):
        half_dim = self.dim // 2
        emb = torch.log(torch.tensor(10000.0)) / (half_dim - 1)
        emb = torch.exp(torch.arange(half_dim, device=t.device) * -emb)
        emb = t[:, None] * emb[None, :]
        emb = torch.cat([torch.sin(emb), torch.cos(emb)], dim=-1)
        return emb
\end{lstlisting}

\subsection{Variance Schedule and Forward Diffusion}

A crucial component of the diffusion process is the variance schedule
$\{\beta_t\}_{t=0}^{T-1}$.  We employ a linear schedule ranging from
$\beta_{\text{start}}=0.0001$ to $\beta_{\text{end}}=0.02$ over
$T=1000$ timesteps.  The scheduler precomputes cumulative products
$\bar{\alpha}_t = \prod_{s=1}^t (1-\beta_s)$ and their square roots for
efficient access during training and sampling.

\begin{lstlisting}[language=Python, caption=Variance Scheduler]
class DDPMScheduler:
    def __init__(self, num_timesteps=1000, beta_start=0.0001, beta_end=0.02, device='cuda'):
        self.num_timesteps = num_timesteps
        self.device = device
        # Linear schedule
        self.betas = torch.linspace(beta_start, beta_end, num_timesteps, device=device)
        self.alphas = 1.0 - self.betas
        self.alphas_cumprod = torch.cumprod(self.alphas, dim=0)
        self.sqrt_alphas_cumprod = torch.sqrt(self.alphas_cumprod)
        self.sqrt_one_minus_alphas_cumprod = torch.sqrt(1.0 - self.alphas_cumprod)

    def perturb_input(self, x0, t, noise=None):
        if noise is None:
            noise = torch.randn_like(x0)
        sqrt_alpha_cumprod_t = self.sqrt_alphas_cumprod[t].view(-1, 1, 1, 1)
        sqrt_one_minus_alpha_cumprod_t = self.sqrt_one_minus_alphas_cumprod[t].view(-1, 1, 1, 1)
        xt = sqrt_alpha_cumprod_t * x0 + sqrt_one_minus_alpha_cumprod_t * noise
        return xt, noise
\end{lstlisting}

\subsection{Training Loop}

The training loop iterates over batches of clean data $x_0$.  For each
batch, a random timestep $t$ is sampled and noise is added to produce
$x_t$.  The network predicts the noise and the mean squared error is
computed.  Gradients are backpropagated and the optimiser updates the
parameters.

\begin{lstlisting}[language=Python, caption=Training Step]
def train_step(model, scheduler, optimizer, x0):
    model.train()
    batch_size = x0.shape[0]
    # Sample random timesteps
    t = torch.randint(0, scheduler.num_timesteps, (batch_size,), device=x0.device)
    # Forward diffusion
    xt, noise = scheduler.perturb_input(x0, t)
    # Predict noise
    noise_pred = model(xt, t)
    # Compute loss
    loss = F.mse_loss(noise_pred, noise)
    optimizer.zero_grad()
    loss.backward()
    optimizer.step()
    return loss.item()
\end{lstlisting}

This step is repeated for the entire training dataset across multiple
epochs.  Typically one trains for 50 epochs with a batch size of
128, although the precise choice depends on hardware and dataset
size.  Gradient clipping can be employed to stabilise training.

\subsection{Sampling Procedures}

\paragraph{DDPM Sampling.}  Once trained, sampling begins from
Gaussian noise $x_T\sim\mathcal{N}(0,I)$.  The DDPM update rule
progresses from $t=T$ down to $t=1$, injecting a small amount of
noise at each step.  An implementation is shown below.

\begin{lstlisting}[language=Python, caption=DDPM Sampler]
@torch.no_grad()
def sample_ddpm(model, scheduler, batch_size, img_size):
    model.eval()
    x = torch.randn(batch_size, *img_size, device=scheduler.device)
    for t in reversed(range(scheduler.num_timesteps)):
        t_tensor = torch.full((batch_size,), t, device=scheduler.device)
        noise_pred = model(x, t_tensor)
        alpha = scheduler.alphas[t]
        alpha_bar = scheduler.alphas_cumprod[t]
        beta = scheduler.betas[t]
        if t > 0:
            noise = torch.randn_like(x)
        else:
            noise = torch.zeros_like(x)
        x = (1/torch.sqrt(alpha)) * (x - (1-alpha)/torch.sqrt(1-alpha_bar) * noise_pred) + torch.sqrt(beta) * noise
    return x
\end{lstlisting}

\paragraph{DDIM Sampling.}  For faster generation we employ the
DDIM sampler.  It uses a strided list of timesteps and sets the
noise scale to zero, resulting in deterministic sampling.  The
following code demonstrates DDIM sampling with 50 inference steps.

\begin{lstlisting}[language=Python, caption=DDIM Sampler]
class DDIMSampler:
    def __init__(self, model, scheduler, device):
        self.model = model
        self.scheduler = scheduler
        self.device = device

    @torch.no_grad()
    def sample(self, batch_size, num_inference_steps=50, eta=0.0):
        self.model.eval()
        x = torch.randn(batch_size, 3, 32, 32, device=self.device)
        step_ratio = self.scheduler.num_timesteps // num_inference_steps
        timesteps = torch.arange(0, self.scheduler.num_timesteps, step_ratio, device=self.device)
        timesteps = torch.flip(timesteps, [0])
        for i in range(len(timesteps)):
            t = timesteps[i]
            t_tensor = t.expand(batch_size)
            noise_pred = self.model(x, t_tensor)
            alpha_bar_t = self.scheduler.alphas_cumprod[t]
            alpha_bar_prev = self.scheduler.alphas_cumprod[timesteps[i-1]] if i > 0 else torch.ones_like(alpha_bar_t)
            pred_x0 = (x - torch.sqrt(1 - alpha_bar_t) * noise_pred) / torch.sqrt(alpha_bar_t)
            dir_xt = torch.sqrt(1 - alpha_bar_prev - eta**2 * (1 - alpha_bar_t / alpha_bar_prev)) * noise_pred
            if eta > 0 and i > 0:
                noise = torch.randn_like(x)
            else:
                noise = torch.zeros_like(x)
            x = torch.sqrt(alpha_bar_prev) * pred_x0 + dir_xt + eta * noise
        return x
\end{lstlisting}

\subsection{Parameter Choices and Results}

Our DDPM implementation uses the parameters listed in
Table~\ref{tab:ddpmparams}.  These settings were found to provide a
balance between training time and sample quality.

\begin{table}[H]
  \centering
  \begin{tabular}{lc}
    \toprule
    \textbf{Parameter} & \textbf{Value} \\
    \midrule
    Image Size & $32\times 32$ \\
    Channels & 3 \\
    Timesteps $T$ & 1000 \\
    $\beta_{\text{start}}$ & 0.0001 \\
    $\beta_{\text{end}}$ & 0.02 \\
    Learning Rate & $10^{-4}$ \\
    Batch Size & 128 \\
    Optimiser & AdamW \\
    Training Epochs & 50 \\
    Base Channels & 64 \\
    Time Embedding Dim & 256 \\
    \bottomrule
  \end{tabular}
  \caption{DDPM training parameters.  Other parameters (e.g. weight
    decay) follow standard defaults.}
  \label{tab:ddpmparams}
\end{table}

\paragraph{Visual Results.}  Figure~\ref{fig:ddpmvsddim} compares
samples generated by DDPM using the full 1000 steps and DDIM using
50 steps.  Despite the drastic reduction in steps, the deterministic
samples of DDIM closely match the quality of the stochastic DDPM
samples.

\begin{figure}[H]
  \centering
  % Placeholders for actual images; in practice include figures via
  % \includegraphics or \maybeincludegraphics commands.  Here we
  % describe the content: the left panel shows a grid of DDPM samples
  % and the right panel shows DDIM samples.
  \fbox{\parbox{0.45\linewidth}{\centering DDPM Samples (1000 steps)\\[2ex]\textit{Insert figure here}}}
  \hfill
  \fbox{\parbox{0.45\linewidth}{\centering DDIM Samples (50 steps)\\[2ex]\textit{Insert figure here}}}
  \caption{Comparison of DDPM and DDIM sampling results.}
  \label{fig:ddpmvsddim}
\end{figure}

\subsection{Summary}

In this chapter we have implemented the key components of diffusion
models: the U‑Net network, the variance schedule, the forward
perturbation and the reverse sampling procedures.  We have shown how
simplified MSE training leads to stable optimisation and how DDIM
sampling accelerates generation.  These implementations form the
foundation for further extensions such as stable diffusion and flow
matching.
